\documentclass[11pt,a4paper]{article}
\usepackage[margin=24mm]{geometry}
\usepackage{amsmath,amssymb,amsthm}
\setlength{\emergencystretch}{2em}
\newtheorem{theorem}{Theorem}
\newtheorem{lemma}[theorem]{Lemma}
\newcommand{\gen}[1]{\langle#1\rangle}
\title{The small diagonal suffixes in the equality case}
\author{Human Compression project: L044 continuation}
\date{}
\begin{document}\maketitle
Retain the complete critical-head hypotheses, the actual intersection
property and all intransitive ambient omissions. Use the explicit
two-prefix-orbit equality interface at a unique interior cut $j=g+1$:
\[
 P=P_{g-1},\quad \operatorname{supp}(P)=K\mathbin{\dot\cup}E,
 \quad |K|+|E|=2g+1,\quad
 \operatorname{Fix}(P)=A\cup\{b\}.
\]
Here $K$ is the head orbit, $E$ is one exterior orbit, $C=C_g$
is the head-side omission orbit and $A=\Omega\setminus C$.
Every label $3,\ldots,g-1$ is paired in the full ambient sense.
The unique crossing of $s=\rho_{g+1}$ is $(b\ a_0)$.
These hypotheses are explicit; this note does not infer them for
other cuts or other numbers of prefix orbits.

\begin{lemma}[The crossing fixes the old prefix support]\label{pure}
The generator $s$ fixes $K\cup E$ pointwise. Consequently it has
the form $s=(b\ a_0)w$, where $w$ is an odd involution on
$A\setminus\{a_0\}$.
\end{lemma}
\begin{proof}
Far commutation makes $s$ centralize $P$. The pure alternating
kernel on $K$ forces it to fix $K$ pointwise, and it preserves
the only other nontrivial $P$-orbit $E$. Thus
$s=(b\ a_0)fw$ with disjoint factors on $E$ and
$A\setminus\{a_0\}$. Suppose $f\ne1$. A nonidentity element
centralizing a transitive permutation group cannot fix a point,
so $f$ is a fixed-point-free involution on $E$. Write $|E|=2m$.

The exact actual prefix intersection identifies the two nontrivial
orbits of $H=G_{3,\ldots,g-1}$ as $T=K\setminus
(D\setminus\{6\})$ and $E$, where $|T|=|K|-5$.
Let $L$ consist of the labels in $I=\{3,\ldots,g-1\}$
having an ambient-eligible edge in $E$, and put $F=I\setminus L$.
No eligible edge can be one of the pairs of $f$, since
$s\in G_i$ for every $i\in I$.

Select one eligible edge in $E$ for each label in $L$ and
project to the $m$ pairs of $f$. These selected edges form
a forest. A cycle would supply, for one selected $i$-edge,
an avoiding path between its two pairs. All the generators
in $H$ commute with $f$; the path lifts to $E$, and $s$
supplies any required switches within a pair. This gives
an actual $G_i$ path between the endpoints of an eligible
$i$-edge, a contradiction. Hence $|L|\le m-1$.

Every eligible edge for a label in $F$ lies in $T$, and each
such label is paired. The acquired selected-localization
theorem, applied to the arbitrary subset $F$ of the interval
$I$ and its invariant orbit $T$, gives
$2|F|\le |T|-1=|K|-6$. Therefore
\[
 g-3=|L|+|F|\le \frac{|E|}{2}-1+\frac{|K|-6}{2}
       =\frac{2g-7}{2}<g-3.
\]
Thus $f=1$. Since $s$ is even, $w$ is odd.
No sign assumption on $f$ was used in the contradiction;
this strengthens the preceding exterior-deficit statement.
\end{proof}

For the remaining statements assume that the diagonal suffix
has the actual marking supplied by the equality path argument:
\[
 A=\{a_0,a'_0,\ldots,a_{t-1},a'_{t-1}\},\qquad
 q_i=\rho_{g+1+i}=(a_{i-1}\ a_i)(a'_{i-1}\ a'_i)
 \quad(1\le i<t).
\]
Here $t=r-g-1$, and $u=\rho_g$ acts on $A$ as
$v=\prod_{i=0}^{t-1}(a_i\ a'_i)$.
The actual prefix $P_g$ is transitive on $C$.
It follows that $u$ moves $b$: otherwise every generator
of $P_g$ would fix $b$. Write $c=b^u\in C\setminus\{b\}$.

\begin{theorem}[The six-point suffix is impossible]\label{six}
The above equality case cannot have $t=3$.
\end{theorem}
\begin{proof}
The generator $s$ centralizes
$q_2=(a_1\ a_2)(a'_1\ a'_2)$. Its unique crossing sends
$b$ to $a_0$; since it fixes $C\setminus\{b\}$ by
Lemma~\ref{pure}, it must fix the only remaining fixed point
$a'_0$ of $q_2$. Its action on the four other points is an
odd involution centralizing the displayed double transposition.
Such an involution is exactly one of
$(a_1\ a_2)$ and $(a'_1\ a'_2)$: a commuting permutation
either preserves both two-point orbits, in which case this
is the odd part of their direct product, or exchanges them,
in which case an involution is a product of two transpositions
and is even. Hence
\[
 s=(b\ a_0)(a_1\ a_2)
 \quad\hbox{or}\quad
 s=(b\ a_0)(a'_1\ a'_2).
\]

Consider the actual omission of the last mark $q_2$.
It contains $P_g$, so $C$ is connected in its orbit graph.
The crossing joins $C$ to $a_0$, and the three edges of
$u|_A$ join each $a_i$ to $a'_i$. The edges of $q_1$
join the pair with index $0$ to the pair with index $1$.
In either displayed form of $s$, its remaining edge joins
the pair with index $1$ to the pair with index $2$.
Thus this omission is transitive on $C\cup A=\Omega$,
contrary to hypothesis. This is an actual orbit argument;
no intersection property is asserted for a projected marking.
\end{proof}

\begin{theorem}[The four-point suffix reduces to one form]\label{four}
If $t=2$, then necessarily
\[
 q_1=(a_0\ a_1)(a'_0\ a'_1),\qquad
 s=(b\ a_0)(a_1\ a'_1),\qquad
 u=(b\ c)(a_0\ a'_0)(a_1\ a'_1)\delta,
\]
where $\delta$ is an odd involution supported on
$C\setminus\{b,c\}$. This statement is a reduction, not
an exclusion of the displayed form.
\end{theorem}
\begin{proof}
By Lemma~\ref{pure}, the restriction of $s$ to
$\{a'_0,a_1,a'_1\}$ is an odd involution, hence a single
transposition. If it is $(a'_0\ a_1)$ or $(a'_0\ a'_1)$,
the crossing, the two pair edges of $u$ and this transposition
join every point of $A$ to $C$ without using $q_1$.
Together with the transitive $P_g$ action on $C$, they make
the last ambient omission transitive. Only
$(a_1\ a'_1)$ remains. The involution $u$ preserves $A$
and $C$, has the stated action on $A$ and exchanges $b,c$.
Its other factor is supported on $C\setminus\{b,c\}$.
The three displayed transpositions are odd, so global evenness
forces $\delta$ to be odd.
\end{proof}

\paragraph{Attribution and scope.}
The first lemma strengthens the exterior-deficit argument using
the same acquired selected-localization theorem and its full
predecessors. The six-point omission proof is the orbit argument
already used in the accepted small anchored-imprimitive proof,
applied here inside the actual ambient group. The literal
four-point reduction records the remaining decorated dihedral
case; its exterior factor has not been removed. The submitted
diagonal row-swap theorem separately excludes the large remaining
centralizer action for $t\ge4$. Other prefix-orbit cases,
endpoint cuts, opposite heads and terminal singleton remain.
No new computation, formal theorem, whole-original or question
credit is claimed.
\end{document}
